\chapter{Sinteză: un singur mecanism, multe aplicații}
\label{ch:sinteza}

Dincolo de ECG (Capitolul~\ref{ch:ecg}), EEG (Capitolul~\ref{ch:eeg}) și
recunoașterea muzicii (Capitolul~\ref{ch:shazam}), aceeași transformată apare în
audio (compresie, anularea zgomotului în căști), în securitate cibernetică
(steganografie și \emph{watermarking}, mesaje ascunse în coeficienții HH,
invizibili ochiului), în finanțe (analiza volatilității pe scări de timp),
în seismologie (separarea undelor P și S) și în astronomie (analiza
semnalelor cosmice slabe, îngropate în zgomot).

Privit dintr-o parte, mecanismul e o alegere de dicționar. O bază wavelet e
un dicționar fix, construit dinainte din condiții matematice, în care
semnalele netede pe porțiuni au reprezentări rare, adică puțini coeficienți
mari și restul aproape de zero. Toate capitolele de mai sus exploatează
exact această raritate, fie tăind coeficienții mici ca zgomot, fie
cheltuind biți doar pe cei mari. Întrebarea firească de aici încolo e dacă
dicționarul trebuie ales dinainte sau poate fi învățat din datele înseși,
ceea ce duce la reprezentări rare cu dicționare antrenate, unde baza se
învață din exemplele pe care urmează să le comprime, în locul unei
construcții pornite din condiții teoretice.

Ce leagă toate aplicațiile din acest document e un singur gest, repetat:
descompune semnalul pe scări (aceeași bancă de filtre din
Secțiunea~\ref{sec:mallat}, aplicată recursiv), taie ce nu contează
(zgomot la denoising, biți la compresie, benzi la analiză) și
reconstruiește cu filtrele de sinteză, exact sau controlat de aproape.
Se schimbă doar ce înseamnă „nu contează”: zgomotul muscular la ECG
(Capitolul~\ref{ch:denoising}), biții pe care ochiul nu-i vede la JPEG2000
(Capitolul~\ref{ch:compresie}), tot ce nu e vârf spectral la recunoașterea
muzicii.

\begin{figure}[htbp]
\centering
\includegraphics[width=\textwidth]{48-synthesis.png}
\caption{Ecranul de sinteză al turului: descompune / taie / reconstruiește,
  aceleași trei verbe explicând ECG, EEG, recunoașterea muzicii, denoising
  și JPEG2000.}
\label{fig:synthesis}
\end{figure}

Localizarea simultană în timp și frecvență, motivul întregului
Capitolul~\ref{ch:wavelets}, e motivul concret pentru care toate aplicațiile
de mai sus funcționează.
